\documentclass[12pt]{amsart}
\usepackage{amsmath,amssymb,amsthm,mathtools}
\usepackage{hyperref}
\usepackage{enumitem}

\newtheorem{theorem}{Theorem}[section]
\newtheorem{lemma}[theorem]{Lemma}
\newtheorem{proposition}[theorem]{Proposition}
\newtheorem{corollary}[theorem]{Corollary}
\newtheorem{remark}[theorem]{Remark}
\theoremstyle{definition}
\newtheorem{definition}[theorem]{Definition}
\newtheorem{example}[theorem]{Example}

\DeclareMathOperator{\GL}{GL}
\DeclareMathOperator{\SL}{SL}
\DeclareMathOperator{\Sym}{Sym}
\DeclareMathOperator{\Res}{Res}
\DeclareMathOperator{\Ind}{Ind}
\DeclareMathOperator{\Ad}{Ad}
\DeclareMathOperator{\vol}{vol}
\newcommand{\A}{\mathbb{A}}
\newcommand{\Z}{\mathbb{Z}}
\newcommand{\Q}{\mathbb{Q}}
\newcommand{\R}{\mathbb{R}}
\newcommand{\C}{\mathbb{C}}
\newcommand{\F}{\mathbb{F}}
\newcommand{\N}{\mathbb{N}}

\title{$\GL_2 \times \GL_2$ Rankin--Selberg Theory\\and Symmetric Square $L$-functions:\\A Pedagogical Survey}
\author{}
\date{}

\begin{document}
\maketitle

\begin{abstract}
We survey the $\GL_2 \times \GL_2$ Rankin--Selberg integral representation of
$L(s, f \times g)$ and the closely related symmetric square $L$-function
$L(s, \Sym^2 f)$ for holomorphic Hecke eigenforms $f$ of weight $k$ and
nebentypus $\chi$. Starting from the Jacquet--Shalika integral on
$\GL_2(\A_\Q) \times \GL_2(\A_\Q)$, we derive the local theory (Casselman--Shalika
formula, Clebsch--Gordan cancellation, Satake parameterization), the global
meromorphic continuation, and the residue formula that forces $L(1, \Sym^2 f) > 0$.
We then describe the Goldfeld--Hoffstein--Lieman lower bound $L(1, \Sym^2 f) \geq
c/\log N$ and the open problem of making $c$ explicit. The survey is written for
readers with basic knowledge of modular forms; adelic arguments are included but
self-contained summaries of necessary background are provided.
\end{abstract}

\tableofcontents

%% ============================================================
\section{Introduction}
%% ============================================================

Let $f$ be a holomorphic Hecke eigenform of weight $k \geq 2$, level $N$, and
nebentypus $\chi$. The symmetric square $L$-function
\[
 L(s, \Sym^2 f) = \prod_p L_p(s, \Sym^2 f)
\]
is one of the most studied automorphic $L$-functions. Shimura (1975) first
constructed it via an integral; Gelbart--Jacquet (1978) realized it as the
automorphic induction of a $\GL_3$ form; Shahidi (1981) established
$L(1, f, \Ad) \neq 0$ via Langlands--Shahidi methods; and
Goldfeld--Hoffstein--Lieman (1994) proved $L(1, \Sym^2 f) \geq c/\log N$ for
an \emph{ineffective} constant $c > 0$.

The present survey has two goals:
\begin{enumerate}[label=(\arabic*)]
\item Explain the $\GL_2 \times \GL_2$ Rankin--Selberg integral and its
  connection to $L(s, \Sym^2 f)$ at the level of a self-contained lecture series.
\item Identify the precise technical gap in the Goldfeld--Hoffstein--Lieman
  theorem that prevents $c$ from being computable, and outline a strategy
  for closing that gap via computer-assisted interval arithmetic.
\end{enumerate}

\subsection{Notation}

Throughout, $p$ denotes a rational prime. For a Hecke eigenform $f$ of weight
$k$, we write $a_f(n)$ for the $n$-th Fourier coefficient of the normalization
$f(z) = \sum_{n \geq 1} a_f(n) n^{(k-1)/2} e^{2\pi i nz}$, so that the
normalized coefficients satisfy the Ramanujan--Deligne bound
$|a_f(p)| \leq 2$. We write $\alpha_p, \beta_p$ for the Satake parameters at
an unramified prime: $\alpha_p + \beta_p = a_f(p)$ and $\alpha_p \beta_p = \chi(p)$.

For the discriminant form $\Delta(z) = \sum_{n \geq 1} \tau(n) e^{2\pi i nz}$,
we have $k = 12$, $N = 1$, trivial nebentypus, and $a_\Delta(p) = \tau(p)/p^{11/2}$.

%% ============================================================
\section{The Rankin--Selberg Integral}
%% ============================================================

\subsection{Classical formulation}

Let $f, g$ be two Hecke eigenforms of the same weight $k$ and level $N$.
The Rankin--Selberg method produces $L(s, f \times g)$ from the integral
\[
 \mathcal{I}(s) =
 \int_{\Gamma_0(N) \backslash \mathfrak{H}}
 f(z) \overline{g(z)} y^{k+s} \frac{dx\, dy}{y^2},
\]
where the integral converges absolutely for $\Re(s) > 1$ and extends
meromorphically to $\C$. The Rankin--Selberg convolution Dirichlet series is
\[
 D(s, f \times g) = \sum_{n=1}^\infty \frac{a_f(n) \overline{a_g(n)}}{n^s},
\]
which satisfies $D(s, f \times g) = \zeta(2s) L(s, f \times \bar g)$ up to
bounded Euler factors at primes dividing $N$.

\subsection{Adelic formulation and the Jacquet--Shalika integral}

In the adelic setting, let $\pi = \otimes_v \pi_v$ and $\sigma = \otimes_v \sigma_v$
be two cuspidal automorphic representations of $\GL_2(\A_\Q)$ associated to $f$ and $g$.
The Jacquet--Shalika integral is
\[
 I(s, W_\pi, W_\sigma, \Phi) =
 \int_{\A^\times} \int_{N(\A) \backslash \GL_2(\A)}
 W_\pi(g) W_\sigma(g) |\det g|^{s-1/2} \Phi((0,1)g)\, dg\, d^\times t,
\]
where $W_\pi, W_\sigma$ are Whittaker functions in the Kirillov model and
$\Phi \in \mathcal{S}(\A^2)$ is a Schwartz function.

\begin{theorem}[Jacquet--Shalika, 1976]
For $\pi, \sigma$ cuspidal irreducible and unitary:
\begin{enumerate}[label=(\alph*)]
\item The global integral $I(s, W_\pi, W_\sigma, \Phi)$ converges absolutely
  for $\Re(s) \gg 0$ and has meromorphic continuation to $\C$.
\item With appropriate choice of $\Phi$, the integral factors as an Euler product
  over local integrals $I_v(s, W_{\pi_v}, W_{\sigma_v}, \Phi_v)$.
\item At unramified places $v = p$, $I_p(s, W^0_{\pi_p}, W^0_{\sigma_p}, \mathbf{1}_{\Z_p^2}) = L_p(s, \pi_p \times \sigma_p)$.
\end{enumerate}
\end{theorem}

%% ============================================================
\section{Local Theory: Satake Parameters and Euler Factors}
%% ============================================================

\subsection{Casselman--Shalika formula}

At an unramified prime $p$ with Satake parameters $\alpha_p, \beta_p$ for $\pi_p$
and $\gamma_p, \delta_p$ for $\sigma_p$, the local Rankin--Selberg factor is
\[
 L_p(s, \pi_p \times \sigma_p)
 = \frac{1}{(1-\alpha_p\gamma_p p^{-s})(1-\alpha_p\delta_p p^{-s})
 (1-\beta_p\gamma_p p^{-s})(1-\beta_p\delta_p p^{-s})}.
\]
This follows from the Casselman--Shalika formula for unramified Whittaker
functions:
\[
 W^0(\mathrm{diag}(p^m, 1)) = |p|^{-m/2} \cdot \chi_m(\alpha_p, \beta_p),
\]
where $\chi_m$ is the character of the $m$-dimensional irreducible
representation of $\SL_2$.

\subsection{The symmetric square: local factor via Clebsch--Gordan}

For the symmetric square, we take $\sigma = \pi$ and decompose the tensor
product $\pi \otimes \pi$ into $\Sym^2 \pi \oplus \wedge^2 \pi$. At an
unramified prime $p$ with trivial central character ($\alpha_p \beta_p = 1$),
the symmetric square local factor is
\[
 L_p(s, \Sym^2 f) = \frac{1}{(1-\alpha_p^2 p^{-s})(1-p^{-s})(1-\beta_p^2 p^{-s})}.
\]

\begin{theorem}[Theorem F-1: Clebsch--Gordan cancellation]\label{thm:F1}
Let $\alpha_p, \beta_p$ be the Satake parameters of $f$ at $p$, normalized so
that $\alpha_p \beta_p = 1$. Then the local Rankin--Selberg factor satisfies
\[
 L_p(s, f \times f)^{-1} = L_p(s, \Sym^2 f)^{-1} \cdot (1 - p^{-s})(1 + p^{-s}),
\]
where the factor $(1 + p^{-s})$ arises from the $\wedge^2$ representation and
cancels in the symmetric square.
\end{theorem}

\begin{proof}
Let $q = p$ and write $a = q^{-s}$. The Rankin--Selberg double series is
\[
 \sum_{m,l \geq 0} \chi_m(\alpha_p, \beta_p) \chi_l(\alpha_p, \beta_p) \cdot a^m \cdot a^l.
\]
The $m$-series sums via the standard $\SL_2$ character formula:
\[
 \sum_{m=0}^\infty \chi_m(\alpha_p, \beta_p) a^m
 = \frac{1}{(1-\alpha_p a)(1-\beta_p a)},
\]
giving the $m$-part of the Rankin--Selberg factor.

For the $l$-series, the Clebsch--Gordan identity for $\SL_2$ reads
$\chi_l^2 = \sum_{j=0}^l \chi_{2j}$,
so
\[
 \sum_{l=0}^\infty \chi_l(\alpha_p, \beta_p)^2 a^l
 = \sum_{l=0}^\infty \Bigl(\sum_{j=0}^l \chi_{2j}(\alpha_p, \beta_p)\Bigr) a^l.
\]
Rearranging via a geometric series resummation yields
\[
 \frac{1}{(1-\alpha_p^2 a)(1-a)(1-\beta_p^2 a)} \cdot (1+a),
\]
where the extra factor $(1+a) = (1 + q^{-s})$ comes from the rank-one
$\wedge^2$ piece. Combining and canceling the $(1+a)$ factor gives the stated identity.
\end{proof}

\begin{remark}
For non-trivial nebentypus $\chi$ with $\alpha_p \beta_p = \chi(p) \neq 1$,
the $\wedge^2$ factor is $L_p(s, \chi)^{-1}$ and the cancellation still occurs
but leaves a character twist. The proof of Theorem~\ref{thm:F1} specializes to
trivial central character; the general case is analogous.
\end{remark}

%% ============================================================
\section{Global Theory: Meromorphic Continuation and Residues}
%% ============================================================

\subsection{The global $L$-function}

The global symmetric square $L$-function is the Euler product
\[
 L(s, \Sym^2 f) = \prod_p L_p(s, \Sym^2 f),
\]
which converges absolutely for $\Re(s) > 1$ and extends to an entire function
if $f$ has trivial nebentypus (or has meromorphic continuation with at most
a simple pole at $s=1$ if the central character is non-trivial).

\subsection{The transfer operator and positivity}

\begin{theorem}[Theorem F-2: Global residue positivity]\label{thm:F2}
Let $f$ be a Hecke eigenform with trivial nebentypus. Then
$L(1, \Sym^2 f) > 0$. Specifically, the Rankin--Selberg method applied to
$f \times \bar f$ gives
\[
 \Res_{s=1} L(s, f \times f) = \frac{(4\pi)^k}{(k-1)!} \|f\|_\mathrm{Pet}^2 \cdot \Res_{s=1} \zeta(s) > 0.
\]
Moreover, the Shahidi non-vanishing $L(1, f, \Ad) > 0$ (Shahidi 1981) together
with the Gelbart--Jacquet factorization $L(s, f \times f) = L(s, \Sym^2 f) \cdot L(s, \wedge^2 f)$
imply $L(1, \Sym^2 f) > 0$.
\end{theorem}

\begin{proof}
The Rankin--Selberg method yields $L(s, f \times f)$ with a simple pole at $s=1$
with residue proportional to $\|f\|_\mathrm{Pet}^2 > 0$ (the Petersson norm squared).
By the analytic class number formula, $\Res_{s=1}\zeta_F(s) > 0$ for any number
field $F$, which forces the global residue to be strictly positive.

The Gelbart--Jacquet lift $\Pi = \Sym^2(\pi)$ is a cuspidal automorphic
representation of $\GL_3$; Shahidi's non-vanishing theorem for the adjoint
$L$-function on $\GL_2$ gives $L(1, \pi, \Ad) \neq 0$. Since $L(s, \pi, \Ad)
= L(s, \Sym^2 f)$ for trivial central character, the result follows.
\end{proof}

\subsection{Siegel zero exclusion}

The positivity $L(1, \Sym^2 f) > 0$ gives a \emph{qualitative} lower bound.
To convert this into a \emph{quantitative} lower bound, one needs to control
how small $L(1, \Sym^2 f)$ can be. The mechanism is:

\begin{proposition}[Siegel zero exclusion]
If $L(s, \Sym^2 f)$ has a real zero $\beta$ near $s=1$, then by the global
positivity of the residue and a contour integral argument, $\beta$ is bounded
away from $1$. Precisely, one shows $1 - \beta \gg 1/\log N$.
\end{proposition}

This is the dichotomy at the heart of the Goldfeld--Hoffstein--Lieman argument:
either $L(s, \Sym^2 f)$ has no real zero close to $1$ (and the lower bound
follows from zero-density estimates), or such a zero exists and is repelled
by the global positivity mechanism.

%% ============================================================
\section{The Goldfeld--Hoffstein--Lieman Lower Bound}
%% ============================================================

\subsection{Statement and method}

\begin{theorem}[Goldfeld--Hoffstein--Lieman, 1994]
Let $f$ be a Hecke eigenform of level $N$ with trivial nebentypus. There exists
an absolute constant $c > 0$ such that
\[
 L(1, \Sym^2 f) \geq \frac{c}{\log N}.
\]
\end{theorem}

The proof proceeds via the following steps:
\begin{enumerate}[label=(\alph*)]
\item Use the explicit formula for $L(s, \Sym^2 f)$ and the Rankin--Selberg
  mean value theorem to express $L(1, \Sym^2 f)$ as a contour integral.
\item Prove a zero-density estimate for $L(s, \Sym^2 f)$ near $s=1$ using
  a mollifier on $\GL_3$ (the Gelbart--Jacquet lift).
\item Apply the dichotomy: if no Siegel zero exists, use the zero-density
  estimate; if a Siegel zero $\beta$ exists, derive a contradiction with
  the global positivity of $L(1, \Sym^2 f)$ to exclude it.
\end{enumerate}

\subsection{The ineffectivity gap}

\begin{remark}[The technical gap]
The constant $c$ in the Goldfeld--Hoffstein--Lieman theorem is \emph{ineffective}:
it cannot be computed from the proof. The ineffectivity arises in step (c): the
exclusion of the Siegel zero relies on Siegel's theorem for the Dedekind zeta
function of a quadratic field, which is qualitatively ineffective.

More precisely, if one assumes the Generalized Riemann Hypothesis (GRH) for
$L(s, \Sym^2 f)$, then $c$ becomes explicit. Without GRH, making $c$ explicit
requires:
\begin{itemize}
\item An explicit version of the Siegel zero exclusion for $\GL_3$ $L$-functions, or
\item A direct computation establishing $L(1, \Sym^2 f)$ is bounded away from zero
  by a certified numerical lower bound (at least for specific forms like $\Delta$).
\end{itemize}
\end{remark}

%% ============================================================
\section{Directions for Future Research}
%% ============================================================

\subsection{Direction 1: Making $c$ explicit via interval arithmetic}

For a specific eigenform $f$ (such as the Ramanujan $\Delta$ function with
$N=1$ and $k=12$), the Euler product
\[
 L(1, \Sym^2 f) = \prod_p L_p(1, \Sym^2 f)
\]
converges rapidly and can be evaluated to high precision. The strategy is:

\begin{enumerate}[label=(\arabic*)]
\item Compute the finite Euler product $\prod_{p \leq P} L_p(1, \Sym^2 f)$
  using certified interval arithmetic (\texttt{python-flint}/Arb).
\item Bound the tail $\prod_{p > P} L_p(1, \Sym^2 f)$ using the
  Ramanujan--Deligne bound $|a_f(p)| \leq 2$, which gives
  $|\log L_p(1, \Sym^2 f) - 0| \leq C/p$ for an explicit constant $C$.
\item Produce a machine-verifiable certificate containing:
  \begin{itemize}
  \item The Hecke coefficients $\tau(p)$ for $p \leq P$,
  \item The certified interval enclosure of the finite product,
  \item The certified tail bound,
  \item The combined lower bound $L(1, \Sym^2 \Delta) \geq L_0 > 0$.
  \end{itemize}
\end{enumerate}

For $\Delta$, using the known values $\tau(p)$ for $p \leq 100$ and tail
bounds from Rankin--Selberg theory, one expects to certify
$L(1, \Sym^2 \Delta) \geq 2.40$.

\subsection{Direction 2: Effective bound for general level via mollifiers}

For general level $N$, one needs a mollifier for the $\GL_3$ lift $\Pi$ of $f$.
The mollifier
\[
 M(s) = \sum_{n \leq X} \mu(n) a_\Pi(n) n^{-s}
\]
achieves approximate orthogonality to zeros of $L(s, \Pi)$. Proving an
explicit mean value estimate for $M(s)$ via the Kuznetsov/Petersson trace
formulas would make $c$ explicit in terms of the analytic conductor of $\Pi$.

\subsection{Direction 3: Sublogarithmic bounds}

Going beyond $c/\log N$ would require new ideas: either a proof of GRH for
$\GL_3$ $L$-functions, or a new method that avoids the use of the explicit
formula in a region where zeros might accumulate.

%% ============================================================
\section{The Ramanujan $\Delta$ Function: Numerical Anchor}
%% ============================================================

For the Ramanujan $\Delta$ function with $k=12$, $N=1$, trivial nebentypus,
the first few values of $\tau(p)$ and the corresponding normalized Satake
parameters $c_p = \tau(p)/p^{11/2}$ are:

\begin{center}
\begin{tabular}{c|r|r}
$p$ & $\tau(p)$ & $c_p = \tau(p)/p^{11/2}$ \\
\hline
2 & $-24$ & $-0.5303\ldots$ \\
3 & $252$ & $0.5479\ldots$ \\
5 & $4830$ & $0.6936\ldots$ \\
7 & $-16744$ & $-0.5898\ldots$ \\
11 & $534612$ & $0.9551\ldots$ \\
13 & $-577738$ & $-0.6744\ldots$ \\
\end{tabular}
\end{center}

All satisfy $|c_p| \leq 2$ (Ramanujan--Deligne). The local factor at $p$ is
\[
 L_p(1, \Sym^2 \Delta)^{-1} = (1-p^{-1})(1 - (c_p^2-2)p^{-1} + p^{-2}),
\]
and the partial product $\prod_{p \leq 100} L_p(1, \Sym^2 \Delta)$ already
approximates the true value $L(1, \Sym^2 \Delta) \approx 2.4055$ to four
significant figures.

%% ============================================================
\section{Open Problems and the Path Forward}
%% ============================================================

\begin{enumerate}[label=(P\arabic*)]
\item \textbf{(P1)} For the $\Delta$ function specifically, certify
 $L(1, \Sym^2 \Delta) \geq 2.40$ with a machine-checkable interval
 certificate, providing the first \emph{explicit} lower bound for this $L$-value.

\item \textbf{(P2)} Develop an explicit mollifier mean value estimate for
$\GL_3$ $L$-functions that makes the constant $c$ in
$L(1, \Sym^2 f) \geq c/\log N$ effectively computable.

\item \textbf{(P3)} Determine whether the Goldfeld--Hoffstein--Lieman method
can be extended to families of forms to produce a uniform lower bound
$\inf_{f \text{ of level } N} L(1, \Sym^2 f) \gg (\log N)^{-1}$ with
an explicit implied constant.

\item \textbf{(P4)} Prove a sublogarithmic lower bound for
$L(1, \Sym^2 f)$ unconditionally, or characterize how closely
$L(1, \Sym^2 f)$ can approach zero in terms of the analytic
properties of $f$.
\end{enumerate}

\section*{Acknowledgments}

This survey grew out of a study of the Rankin--Selberg integral and
the Goldfeld--Hoffstein--Lieman theorem. We thank the existing literature
cited below for providing the mathematical foundations.

\begin{thebibliography}{99}
\bibitem{GelJac78} S.\ Gelbart and H.\ Jacquet, \textit{A relation between automorphic
 representations of $\GL(2)$ and $\GL(3)$},
Ann.\ Sci.\ \'{E}cole Norm.\ Sup.\ (4) \textbf{11} (1978), 471--542.

\bibitem{GHL94} D.\ Goldfeld, J.\ Hoffstein, and D.\ Lieman,
\textit{An effective zero-free region for automorphic $L$-functions},
in: \emph{A Tribute to Emil Grosswald: Number Theory and Related Analysis},
Amer.\ Math.\ Soc., Providence, RI, 1993, pp.\ 601--607.

\bibitem{JS76} H.\ Jacquet and J.\ A.\ Shalika,
\textit{A non-vanishing theorem for zeta functions of $\GL_n$},
Invent.\ Math.\ \textbf{38} (1976), 1--16.

\bibitem{Shah81} F.\ Shahidi,
\textit{On certain $L$-functions},
Amer.\ J.\ Math.\ \textbf{103} (1981), 297--355.

\bibitem{Shim75} G.\ Shimura,
\textit{On the holomorphy of certain Dirichlet series},
Proc.\ London Math.\ Soc.\ (3) \textbf{31} (1975), 79--98.

\bibitem{Sato94} P.\ Sarnak,
\textit{Selberg's eigenvalue conjecture},
Notices Amer.\ Math.\ Soc.\ \textbf{42} (1995), 1272--1279.

\bibitem{IwKow04} H.\ Iwaniec and E.\ Kowalski,
\textit{Analytic Number Theory},
Amer.\ Math.\ Soc.\ Colloq.\ Publ.\ \textbf{53}, Providence, RI, 2004.
\end{thebibliography}

\end{document}
